%% The following is a directive for TeXShop to indicate the main file
%%!TEX root = diss.tex

\chapter{Abstract}

The rise in the amount of graph data that are being generated has prompted an increased interest in storing and processing these datasets.
At this point, there has been nearly a decade of research effort toward developing efficient graph processing systems.
The standard processing pipeline first ingests the data into an in-memory representation before running any processing tasks.
If the graph dataset is too large to fit in memory on a single node, the graph is partitioned and distributed so that it fits in aggregate memory on a cluster.
Whenever there are new data to analyze, all this effort is repeated.

One representation does not suit all workloads or graphs, and changing a custom storage engine that is tightly coupled with the graph processing system and its data representation can be difficult and error-prone.
Using a database, in contrast, appears to be the natural choice to manage evolving data.
However, databases remain relatively underutilized for running in-situ analytics \\PM{this is not the correct term for what you are trying to describe.} tasks for various reasons ranging from difficulty to program iterative tasks to poor performance on analytics benchmarks. \marginpar[Trying the note]{trying the margin note}

We take a middle-of-the-road approach and created \project{}, a graph database built on top of a mature key-value store that provides production quality data management while allowing flexibility in prototyping different data layouts easily without having to reimplement the storage stack.
We enable fast in-database analytics by using at rest data layouts that match closely with the in-memory layout needed by the analytics tasks. We also eliminate the need to run ETL tools or preprocess the graph.
\project{} is faster than state-of-the-art commercial graph databases on standard graph benchmarks and can withstand a high insertion rate of X ops/sec.
\project{} achieves upto X\% performance compared to other specialized out-of-core graph processing systems \PM{examples}.

